\section{Shape inequalities for finite step functions}
\label{shape:section}

On a positive constant interval, the curvature formula
\eqref{eq:second-curvature} contains the indefinite expression
$z^2-12vw^2+6v^4$.  The next section will make its expectation positive
by adding a spatial derivative of zero expectation.  Our task here is to
prove the shape inequalities that make this possible: we compare the
backward drift with a logarithmic derivative of the forward density.

We work with a fixed positive finite step coefficient and a smooth
terminal condition.  The estimates in this section may depend on these
fixed data.  Section~\ref{sec:curvature} will pass only the resulting
integrated inequality to general coefficients, using
Lemma~\ref{lem:regularity}.

Fix a partition $0=t_0<t_1<\cdots<t_N=1$, constants
$0<c_1\le\cdots\le c_N\le M$, and
$\gamma(t)=c_i$ on $[t_{i-1},t_i)$.  Let $\Phi$ solve
\[
 \Phi_t=-\frac12(\Phi_{xx}+\gamma\Phi_x^2),
 \qquad \Phi(1,x)=M^{-1}\log(2\cosh(Mx)).
\]
On a step where $\gamma=c$, put
\[
 l=c\Phi,\qquad r=l_x,\qquad v=r_x,
 \qquad w=r_{xx},\qquad z=r_{xxx}.
\]
Spatial derivatives of $\Phi$ agree across step boundaries; those of $l$
are rescaled with $c$.  Within a step, $h=e^l$ satisfies
$h_t=-\tfrac12h_{xx}$.  In backward time $\tau$ this gives
\begin{equation}\label{shape:backward-r}
 r_\tau=\frac12r_{xx}+rr_x.
\end{equation}
Let $X_0=0$ and $\dd X_t=r(t,X_t)\dd t+\dd B_t$.  The drift is bounded
and spatially Lipschitz for these fixed data.  Its density at positive times
will be denoted by $f(t,x)$.  Factor it as
\[
 f=e^l g,\qquad s=-\partial_x\log g.
\]
The factor $e^l$ contains the backward solution.  The remaining factor
$g$ evolves by the forward heat equation within each step.  Indeed, the
transition density from time $a$ to time $t$ in the same step is
\[
 p_{t-a}(x-y)\frac{h(t,x)}{h(a,y)},\qquad
 p_u(x)=(2\pi u)^{-1/2}e^{-x^2/(2u)}.
\]
It integrates to one because $h(a,\cdot)=P_{t-a}h(t,\cdot)$, and
differentiation verifies the forward equation with drift $r=h_x/h$.
Consequently $g_t=\tfrac12g_{xx}$.  On the first step,
$g(t,x)=p_t(x)/h(0,0)$, so $s=x/t$.  At a forward jump
$c_+-c_-=\delta\ge0$, continuity of $f$ gives
\begin{equation}\label{tail:g-jump}
 g_+=g_-e^{-\delta\Phi}.
\end{equation}
In particular $g$ is positive and even, and $s$ is odd.

The comparison we need has three parts.  On the positive half-line the
drift $r$ is increasing and concave; its relative backward-time change
$r_\tau/r$ is nonpositive and nondecreasing in $x$.  The score $s$ is also
increasing and concave, and its ratio to $r$ is nondecreasing in $x$.
Lemmas~\ref{shape:backward} and~\ref{shape:forward} state these claims
precisely.  We first record the boundary estimates needed in their proofs.

\subsection{Boundary estimates and comparison}

A differentiated exponential error means that, for each fixed finite
derivative order, the error and all its derivatives up to that order
are $O(e^{-\kappa x})$ as $x\to+\infty$, for some $\kappa>0$.
Constants may change with the derivative order.

\begin{lemma}[Differentiated tails]\label{tail:class}
On each closed backward step, uniformly in its time coordinate,
\begin{equation}\label{tail:backward}
 l(t,x)=cx+C(t)+O(e^{-\kappa x})\qquad(x\to+\infty)
\end{equation}
with differentiated exponential error.  The functions are even, so the
negative tail is determined by symmetry.

For the forward heat factor $g$, on every closed forward step with
time bounded away from zero,
\begin{equation}\label{tail:forward}
 g(t,x)=\exp\left(-\frac{x^2}{2t}+\beta(t)x+C_0(t)\right)
                 (1+O(e^{-\kappa x}))\qquad(x\to+\infty),
\end{equation}
with differentiated relative error.  In particular,
\begin{equation}\label{tail:score}
 s(t,x)=\frac{x}{t}-\beta(t)+O(e^{-\kappa x}),\qquad
 s_x(t,x)=\frac1t+O(e^{-\kappa x}),\qquad
 s_{xx}(t,x)=O(e^{-\kappa x}).
\end{equation}
The assertions include one-sided values at every step boundary.
\end{lemma}

The proof uses only Gaussian convolution and the finite-step jump rules;
none of the shape inequalities proved below enters it.

\begin{proof}
Here is a Gaussian estimate that also controls differentiation.  Suppose
$q\in C^m(\mathbb R)$ satisfies, for $0\le j\le m$,
\[
 |q^{(j)}(y)|\le A e^{-\kappa y}\quad(y\ge0),\qquad
 |q^{(j)}(y)|\le A e^{K|y|}\quad(y<0).
\]
For $Y_x=\lambda x+b+\sqrt V Z$, where $Z$ is standard normal,
$0<\lambda_*\le\lambda\le\lambda^*$, $|b|\le B$, and $0\le V\le V_*$,
\begin{equation}\label{tail:gaussian-estimate}
 \partial_x^j\mathbb E q(Y_x)
       =\lambda^j\mathbb E q^{(j)}(Y_x)=O(e^{-\kappa' x})
       \qquad(0\le j\le m)
\end{equation}
uniformly in these parameters, for some $\kappa'>0$.
Indeed, write $\mu=\lambda x+b$ and take $x$ large enough that $\mu>0$.
On $Y_x\ge\mu/2$, use $Ae^{-\kappa\mu/2}$.  On its complement,
the global bound $|q^{(j)}(y)|\le A e^{K|y|}$, Cauchy--Schwarz, and
$\mathbb P(Y_x<\mu/2)\le e^{-\mu^2/(8V_*)}$ give a bound
$C\exp(K\mu-\mu^2/(16V_*))$.  This is smaller than an exponential
in $x$; the case $V=0$ is immediate.  Exponential integrability of the
Gaussian justifies differentiation under the expectation.  Adjusting the
constant covers the remaining bounded range of $x$.

For backward evolution, suppose an even starting function $l_0$ has
$l_0(x)=cx+C+\rho(x)$ with differentiated exponential error.  Set
$q(y)=e^{l_0(y)-cy-C}-1$.  This $q$ obeys the hypotheses above: the
positive tail follows by differentiating the exponential, and evenness of
$l_0$ gives at most exponential growth for $q$ and its derivatives on the
negative tail.  Gaussian tilting gives
\[
 P_\tau(e^{l_0})(x)
 =e^{cx+C+c^2\tau/2}
   \left(1+\mathbb E q(x+c\tau+\sqrt\tau Z)\right).
\]
Apply \eqref{tail:gaussian-estimate} uniformly for bounded $\tau$ and take
the logarithm.  Its relative error tends uniformly to zero, so differentiating
$\log(1+\cdot)$ proves \eqref{tail:backward} with all the stated derivatives.
The initial condition starts this induction since, for $x>0$,
\[
 \frac cM\log(2\cosh Mx)
       =cx+\frac cM\log(1+e^{-2Mx}).
\]
At a backward boundary, $l$ is multiplied by $c_-/c_+>0$, which
preserves the required tail class.  Finite induction proves the first claim.

Suppose at a positive time $a$ that
$g(a,y)=e^{-y^2/(2a)+\beta y+C}(1+q(y))$ has the required tail class.
Evenness of $g$ implies the same exponential-growth bound for the
derivatives of $q$ on the negative tail.  Completing the square gives
\begin{align*}
 P_hg(a,\cdot)(x)
 &=\sqrt{\frac a{a+h}}\,
   \exp\left(-\frac{x^2}{2(a+h)}
      +\frac{a\beta}{a+h}x+C+\frac{\beta^2ah}{2(a+h)}\right)\\
 &\quad\times
   \left(1+\mathbb E q\left(
        \frac a{a+h}x+\frac{\beta ah}{a+h}
           +\sqrt{\frac{ah}{a+h}}Z\right)\right).
\end{align*}
On a fixed interval $0\le h\le H$, the coefficient of $x$ in the
Gaussian mean is at least $a/(a+H)>0$.  Thus
\eqref{tail:gaussian-estimate} proves closure, uniformly including $h=0$.
At a jump, write $\Phi=x+D+\rho$ on the positive tail.  Formula
\eqref{tail:g-jump} replaces $(\beta,C)$ by
$(\beta-\delta,C-\delta D)$ and replaces $1+q$ by
$(1+q)e^{-\delta\rho}$, preserving the differentiated relative error.
Finite induction proves \eqref{tail:forward}.  Logarithmic differentiation
then gives \eqref{tail:score}.
\end{proof}



We will use the following elementary form of comparison.  If on a compact
space-time cylinder
\[
 u_t-\tfrac12u_{xx}-b u_x-c u\ge0,\qquad c\le C,
\]
and $u\ge0$ initially while $u\ge-\varepsilon$ on the lateral boundary,
then
\begin{equation}\label{shape:barrier}
 u(t,x)\ge-\varepsilon e^{(\max(C,0)+1)(t-t_0)}.
\end{equation}
To see this, add the positive exponential on the right to $u$ and multiply
by $e^{-(\max(C,0)+1)(t-t_0)}$.  The resulting zeroth-order coefficient
is at most $-1$.  A negative space-time minimum, which cannot be on the
parabolic boundary, has nonpositive time derivative, zero spatial gradient,
and nonnegative second spatial derivative, contradicting its differential
inequality.  The same proof applies in a ball in $\mathbb R^d$ with
$\tfrac12\Delta$ in place of $\tfrac12\partial_{xx}$.
In particular, to exhaust the line or half-line it suffices that the
lateral boundary errors tend uniformly to zero and that $C$ is independent
of the truncation.  No uniform bound on the first-order coefficient $b$
over the exhausting cylinders is needed.

\subsection{Backward shape inequalities}

\begin{lemma}[Backward shapes]\label{shape:backward}
For the fixed step data above, $r$ is odd and $v>0$.  On $x>0$,
\begin{equation}\label{shape:backward-signs}
 w\le0,\qquad
 H:=\frac{r_\tau}{r}=\frac{w}{2r}+v\le0,
 \qquad H_x\ge0.
\end{equation}
The quotient defining $H$ extends smoothly and evenly at zero.
\end{lemma}

\begin{proof}
Evenness is preserved by both heat evolution and rescaling.  If $l_0''>0$,
the logarithm of its heat evolution obeys
\[
 \partial_{xx}\log P_\tau e^{l_0}(x)
 =\mathbb E_\nu l_0''(Y)
     +\operatorname{Var}_\nu(l_0'(Y))>0,
\]
where $\nu$ is the Gaussian law of $Y=x+\sqrt\tau Z$ tilted by
$e^{l_0(Y)}$; at $\tau=0$ use $l_0''(x)>0$ directly.  Since the terminal
second derivative is $cM\operatorname{sech}^2(Mx)>0$ and backward
rescalings are positive, this proves $v>0$.  Thus $r(x)>0$ for $x>0$.
All required derivatives are bounded on closed steps, by smoothness on
compact sets and Lemma~\ref{tail:class} at infinity.

Differentiating \eqref{shape:backward-r} twice gives
\[
 w_\tau=\tfrac12w_{xx}+r w_x+3v w.
\]
The terminal $w=-2cM^2\operatorname{sech}^2(Mx)\tanh(Mx)$ is
nonpositive for $x>0$.  The lateral values are $w(\tau,0)=0$ and
$w(\tau,A)=o_A(1)$ uniformly on a closed step.  Apply
\eqref{shape:barrier} to $-w$, with $C=3\sup v$, and let $A\to\infty$.
Backward jumps multiply $w$ by a positive constant, so the sign holds
throughout.

Because $r(x)=v(0)x+O(x^3)$ with $v(0)>0$, odd smoothness makes
$w/r$ smooth and even at zero.  Direct differentiation of
$r_\tau=rH$ gives
\begin{align}
 H_\tau&=\tfrac12H_{xx}+(v/r+r)H_x+vH,
     \label{shape:H-evolution}\\
 (H_x)_\tau&=\tfrac12(H_x)_{xx}+(v/r+r)(H_x)_x
       +\bigl((v/r)_x+2v\bigr)H_x+wH.
     \label{shape:Hx-evolution}
\end{align}
At the terminal endpoint of the last step,
\[
 H=(cM-M^2)\operatorname{sech}^2(Mx)\le0,
 \qquad H_x\ge0\quad(x>0).
\]
For \eqref{shape:H-evolution}, write
$v/r+r=x^{-1}+\eta(x)$, where $\eta$ is smooth odd and $\eta(x)=O(x)$
near zero.  The radial lift $\widehat H(y)=H(|y|)$ in $\mathbb R^3$
satisfies
\[
 \widehat H_\tau=\tfrac12\Delta\widehat H
       +\eta(|y|)\frac y{|y|}\cdot\nabla\widehat H+v(|y|)\widehat H.
\]
Both the lifted function and vector field are smooth at zero: a smooth
even function lifts smoothly, and $\eta(x)/x$ is smooth even.
Since $H\to0$ uniformly at infinity, comparison on balls followed by
exhaustion gives $H\le0$.

For $H_x$, use intervals $[\varepsilon,A]$.  The source $wH$ is
nonnegative, and
\[
 (v/r)_x=\frac{w}{r}-\frac{v^2}{r^2}<0,
 \qquad (v/r)_x+2v\le2\sup v.
\]
The lateral errors are $H_x(\tau,\varepsilon)=O(\varepsilon)$,
uniformly by smooth evenness, and $H_x(\tau,A)=o_A(1)$ by the tails.
Formula~\eqref{shape:barrier}, with a bound independent of
$\varepsilon,A$, proves $H_x\ge0$ on taking the two limits.

Finally, at a backward jump write $r_{\rm new}=\alpha r_{\rm old}$,
where $0<\alpha\le1$.  Its derivatives rescale identically, and hence
\[
 H_{\rm new}=H_{\rm old}-(1-\alpha)v_{\rm old},\qquad
 (H_x)_{\rm new}=(H_x)_{\rm old}-(1-\alpha)w_{\rm old}.
\]
These preserve both signs.  Induction over the finitely many steps
completes the proof.
\end{proof}

\subsection{Forward score inequalities}

The logarithmic derivative $s=-g_x/g$ measures the part of the density
not contained in $e^l$.  We need both its concavity and a comparison with
$r$, rather than concavity of either function alone.
The density transform and its use on constant intervals also appear in
Lopatto's positive-temperature analysis~\cite[Sections~3, 7, and 9]{Lopatto2026}.
Here we prove the additional ratio sign needed by the coercive identity.

\begin{lemma}[Forward score and ratio]\label{shape:forward}
For every $t>0$, $s$ is odd and, with $a=s_x$,
\begin{equation}\label{shape:forward-signs}
 a\ge\frac1t>0,\qquad s_{xx}\le0\quad(x>0),\qquad
 j:=a-\frac sr v=r\partial_x(s/r)\ge0\quad(x\in\mathbb R).
\end{equation}
The ratios and $j$ extend smoothly at zero.  At every forward step
boundary, $j$ is unchanged.
\end{lemma}

\begin{proof}
Evenness and positivity of $g$ give odd smoothness of $s$.  Within a step,
the heat equation for $g$ gives
\begin{align}
 s_t&=\tfrac12s_{xx}-s s_x,
       \label{shape:s-evolution}\\
 a_t&=\tfrac12a_{xx}-s a_x-a^2,
       \label{shape:a-evolution}\\
 (s_{xx})_t&=\tfrac12(s_{xx})_{xx}
                  -s(s_{xx})_x-3a s_{xx}.
       \label{shape:sxx-evolution}
\end{align}
On the first interval $s=x/t$, so the first two claims hold explicitly.
At a jump of size $\delta=c_+-c_-$, \eqref{tail:g-jump} gives
\begin{equation}\label{shape:s-jump}
 s_+=s_-+\delta\Phi_x,
 \quad a_+=a_-+\delta\Phi_{xx},
 \quad (s_{xx})_+=(s_{xx})_-+\delta\Phi_{xxx}.
\end{equation}
Here $\Phi_{xx}>0$ and $\Phi_{xxx}\le0$ on the positive half-line by
Lemma~\ref{shape:backward}.

For subsequent intervals, put $u=a-1/t$.  Equation
\eqref{shape:a-evolution} becomes the linear equation
\[
 u_t=\tfrac12u_{xx}-s u_x-(a+1/t)u.
\]
For the fixed smooth solution, $a$ is bounded on a closed step; this
follows from \eqref{tail:score} and compactness and does not presume its
sign.  Thus the zeroth-order coefficient is bounded above independently
of a spatial truncation.  Initially $u\ge0$, by the preceding step and
\eqref{shape:s-jump}; at both infinities $u\to0$ uniformly by
Lemma~\ref{tail:class}.  Comparison on $[-A,A]$ and exhaustion prove
$a\ge1/t$.  Next apply comparison to $-s_{xx}$ on $[0,A]$ using
\eqref{shape:sxx-evolution}.  Its zeroth coefficient is $-3a\le0$,
its initial values are nonnegative, and its boundary values tend to
zero, at zero by oddness and at infinity by \eqref{tail:score}.  This
proves $s_{xx}\le0$.

For the ratio inequality, first note that $s(x)>0$ and $r(x)>0$ for
$x>0$.  On this half-line define
\[
 d_s=\frac{s_x}{s},\qquad d_r=\frac{r_x}{r},\qquad \rho=d_s-d_r.
\]
Then $j=s\rho$, so it suffices to prove $\rho\ge0$.
On the first step, $s=x/t$ and the concavity of $r$ gives
$r(x)=\int_0^x v(y)\dd y\ge xv(x)$; hence
$\rho=1/x-v/r\ge0$.
Within any later step, forward time reverses \eqref{shape:backward-r},
so $r_t=-rH$.  Consequently
\[
 (d_s)_t=\tfrac12(d_s)_{xx}+d_s(d_s)_x-s_{xx},
 \qquad (d_r)_t=-H_x.
\]
Also
$\tfrac12(d_r)_{xx}+d_r(d_r)_x
 =\tfrac12(w/r)_x=H_x-w$.
Subtracting gives
\begin{equation}\label{shape:G-evolution}
 \rho_t=\tfrac12\rho_{xx}+d_s\rho_x+(d_r)_x\rho
                         +2H_x-w-s_{xx}.
\end{equation}
The source is nonnegative by the signs already proved, while
$(d_r)_x=w/r-v^2/r^2<0$.

Here are the endpoint details for this last comparison.  On a closed
step away from time zero, $v(t,0)>0$ is continuous and bounded below,
and $a(t,0)\ge1/t>0$.  Odd Taylor expansions therefore give, uniformly,
\[
 d_s=x^{-1}+O(x),\qquad d_r=x^{-1}+O(x),\qquad \rho=O(x)
       \quad(x\downarrow0).
\]
At infinity, \eqref{tail:backward} and \eqref{tail:score} give
$d_r=O(e^{-\kappa x})$ and
$d_s=(x-\beta(t)t)^{-1}+O(e^{-\kappa' x})$, so $\rho\to0$ uniformly.
Use \eqref{shape:barrier} on $[\varepsilon,A]$, with upper
zeroth-order bound zero, and let $\varepsilon\downarrow0$, $A\to\infty$.
This proves $\rho\ge0$ provided it is nonnegative at the start of the step.

To verify that condition, write at a forward jump
\[
 r_+=\alpha r_-,\quad v_+=\alpha v_-,\quad
 s_+=s_-+\lambda r_-,\quad a_+=a_-+\lambda v_-,
 \qquad
 \alpha=\frac{c_+}{c_-}>0,\quad\lambda=\frac{c_+-c_-}{c_-}\ge0.
\]
Then direct cancellation gives
\[
 j_+=a_-+\lambda v_-
     -\frac{s_-+\lambda r_-}{\alpha r_-}\alpha v_-=j_-.
\]
Since $\rho_+=j_+/s_+$ for $x>0$, its initial sign is preserved.
This proves the ratio claim by induction.  Smoothness of $s/r$ at zero
follows by factoring $x$ from both odd functions; in particular $j$ is
even and $j(0)=0$.  Reflection gives the assertion on the whole line.
\end{proof}

\subsection{Variables for the integral certificate}

We now express the shape inequalities as nonnegative quantities suited
to polynomial algebra.  Set
\[
 R=r^2,\quad K=rs,\quad b=-w/r,\quad
 j=s_x-(s/r)v=r(s/r)_x.
\]
The first two quantities are nonnegative by the common signs of $r$ and
$s$, and $b\ge0$ expresses the concavity of $r$.  The ratio comparison
is precisely $j\ge0$.  The remaining three quantities encode the
backward-time and forward-concavity inequalities:
\[
 B=b-2v=-2H,\qquad
 S=z+bv-2Rb=2rH_x,\qquad L=-r s_{xx}.
\]
All are nonnegative.  The following lemma collects their derivative
rules and the integration-by-parts identity that will be used in the
coercive calculation.

\begin{lemma}[Certificate signs and integration by parts]\label{shape:ibp}
At a fixed positive time in a constant step, the functions
$R,K,v,b,j,L,B,S$ defined above are even and extend smoothly across zero.
They satisfy
\[
 R,K,v,j,L,B,S\ge0.
\]
Writing $\mathcal D=r\partial_x$, one has
\begin{equation}\label{shape:D-rules}
 \mathcal D(R,K,v,b,j)
  =\bigl(2Rv,\ 2vK+Rj,\ -Rb,\ -z-bv,\ -jv+Kb-L\bigr).
\end{equation}
For every polynomial $F$ in $R,K,v,b,j$,
\begin{equation}\label{shape:ibp-identity}
 \mathbb E_f\!\left[\mathcal DF+(v+R-K)F\right]=0.
\end{equation}
Every term in this identity is absolutely integrable.
\end{lemma}

\begin{proof}
The signs of $R,K,v,j,L$ follow from Lemmas~\ref{shape:backward} and
\ref{shape:forward}; for $K$ use the common sign of the odd functions
$r,s$.  Since $H=v-b/2$, we have $B=-2H\ge0$.
On $x>0$, differentiation gives
\[
 S=2rH_x\ge0.
\]
All these expressions are even.  The ratios $w/r$ and $s/r$ are smooth
at zero since their numerators and denominator are odd and $r_x(0)>0$.
The inequalities consequently extend there by continuity.
For example $b(0)=-z(0)/v(0)$, $j(0)=0$, and $S(0)=0$.

The derivative rules follow directly from the definitions.  The only
one involving the forward score is
\[
 rj_x=r s_{xx}-a v-sw+\frac{s v^2}{r}
     =-jv+Kb-L,
\]
which gives the last rule in \eqref{shape:D-rules}.

The density has logarithmic derivative $f_x/f=r-s$.  Therefore
\[
 \partial_x(rFf)
   =\left(\mathcal DF+(v+r(r-s))F\right)f
   =\left(\mathcal DF+(v+R-K)F\right)f.
\]
By the differentiated tails, $r$ is bounded and tends to $c$ at the
positive tail, $v,w,z,b$ and their derivatives decay exponentially,
$s=x/t+O(1)$, $j=1/t+O((1+x)e^{-\kappa x})$, and $K=O(1+x)$.
Every polynomial $F$ and its differentiated expression thus have at
most polynomial growth.  Moreover
\[
 f(t,x)\le C\exp\left(-\frac{x^2}{2t}+C|x|\right).
\]
These bounds prove absolute integrability and $rFf\to0$ at both
infinities.  Integrating the displayed derivative on $[-A,A]$ and
letting $A\to\infty$ proves \eqref{shape:ibp-identity}.
\end{proof}

